% BpfReJIT: Separating Mechanism and Policy in Kernel eBPF JIT Optimization
% ACM SIGCONF format, anonymous submission

\documentclass[sigconf,anonymous]{acmart}

\usepackage{microtype}
\usepackage{booktabs}
\usepackage{graphicx}
\usepackage{listings}
\usepackage{xcolor}
\usepackage{multirow}
\usepackage{amsmath}
\usepackage{url}
\usepackage{subcaption}
\usepackage{array}
\usepackage{tabularx}

% ACM copyright / DOI setup (for submission, leave as default)
\setcopyright{acmlicensed}
\copyrightyear{2026}
\acmYear{2026}
\acmDOI{XXXXXXX.XXXXXXX}
\acmConference[OSDI '26]{20th USENIX Symposium on Operating Systems Design and Implementation}{July 2026}{Santa Clara, CA, USA}
\acmBooktitle{20th USENIX Symposium on Operating Systems Design and Implementation (OSDI '26), July 2026, Santa Clara, CA, USA}
\acmPrice{15.00}
\acmISBN{978-1-4503-XXXX-X/26/07}

% Code listing style
\lstset{
  basicstyle=\ttfamily\footnotesize,
  keywordstyle=\color{blue}\bfseries,
  commentstyle=\color{gray},
  stringstyle=\color{orange},
  breaklines=true,
  numbers=left,
  numberstyle=\tiny\color{gray},
  frame=single,
  xleftmargin=1em,
  xrightmargin=0.5em,
}

% Compact table spacing
\setlength{\tabcolsep}{4pt}

\begin{document}

\title{BpfReJIT: Separating Mechanism and Policy\\in Kernel eBPF JIT Optimization}

\author{Anonymous Authors}
\affiliation{\institution{Anonymous Institution}}
\email{anonymous@example.com}

\begin{abstract}
The Linux kernel eBPF JIT compiler uses a fixed, one-size-fits-all code
generation strategy. Comparing its output against LLVM~-O3 on identical BPF
programs reveals a 15\% execution-time gap and 2$\times$ code-size gap, rooted
in three backend lowering deficiencies: byte-load recomposition, absent
conditional moves, and redundant register saves. A natural fix---hardcoding
these optimizations into the kernel---fails because some lowering choices are
\emph{policy-sensitive}: forcing \texttt{cmov} on predictable branches
degrades performance by 34.5\% ($p{=}0.002$), while the same optimization
yields 5$\times$ speedup on unpredictable branches.

We present \textsc{BpfReJIT}, a framework that separates JIT optimization
\emph{mechanism} (kernel-enforced safety) from \emph{policy} (userspace
profitability decisions). Through a new \texttt{BPF\_PROG\_JIT\_RECOMPILE}
syscall, privileged userspace specifies per-site native instruction choices
for already-verified BPF programs; the kernel validates and re-JITs without
re-verification. Four rule families (conditional select, wide load, rotate
fusion, address calculation) cover the identified gaps. On
microbenchmarks, ROTATE achieves 1.24$\times$ speedup with 15\% code
reduction ($p{=}0.002$). On 60 production BPF programs from Calico, Katran,
and Suricata, all recompilations succeed with zero syscall failures. A
declarative pattern extension reduces the cost of adding new optimization
variants from 1{,}018 LOC (kernel+userspace) to 32 LOC (userspace only).
\end{abstract}

\begin{CCSXML}
<ccs2012>
  <concept>
    <concept_id>10011007.10011006.10011008.10011009.10011015</concept_id>
    <concept_desc>Software and its engineering~Compilers</concept_desc>
    <concept_significance>500</concept_significance>
  </concept>
  <concept>
    <concept_id>10011007.10011074.10011099.10011102.10011103</concept_id>
    <concept_desc>Software and its engineering~Software performance</concept_desc>
    <concept_significance>300</concept_significance>
  </concept>
  <concept>
    <concept_id>10010520.10010553.10010562</concept_id>
    <concept_desc>Computer systems organization~Operating systems</concept_desc>
    <concept_significance>300</concept_significance>
  </concept>
</ccs2012>
\end{CCSXML}

\ccsdesc[500]{Software and its engineering~Compilers}
\ccsdesc[300]{Software and its engineering~Software performance}
\ccsdesc[300]{Computer systems organization~Operating systems}

\keywords{eBPF, JIT compilation, kernel optimization, instruction selection, mechanism/policy separation}

\maketitle

%% ============================================================
\section{Introduction}
\label{sec:intro}
%% ============================================================

eBPF has become the dominant in-kernel extension mechanism, powering
Cilium's network policy enforcement~\cite{cilium}, Meta's Katran load
balancer~\cite{katran}, the \texttt{sched\_ext} scheduling
framework~\cite{sched_ext}, and kernel tracing infrastructure. Every
deployed BPF program passes through the kernel's built-in JIT compiler: a
single-pass, linear-scan translator that converts verified BPF bytecode to
native machine code. This compiler was designed for simplicity, correctness,
and fast compilation---not output code quality.

\textbf{How large is the gap?} We present the first systematic
native-code-level comparison of the kernel JIT against llvmbpf, a userspace
LLVM~-O3 JIT that compiles identical BPF bytecode. Across 31 authoritative
microbenchmarks on an Intel Arrow Lake-S platform (10$\times$1000 iterations,
CPU-pinned, performance governor), llvmbpf achieves an execution-time geomean
of \textbf{0.849$\times$} (95\% CI [0.834, 0.865]) and produces native code
that is \textbf{2$\times$ smaller} (code-size geomean 0.496$\times$).
JIT-dump analysis attributes 89\% of the instruction surplus to three
backend deficiencies: byte-load recomposition (50.7\%), absent conditional
moves (19.9\%), and fixed callee-saved registers (18.5\%).

\textbf{Why not just fix the kernel?} A natural response is to add these
optimizations as kernel peepholes. We show that this fails for a fundamental
reason: some lowering choices are \emph{policy-sensitive}---the same
optimization helps one workload and hurts another. Forcing \texttt{cmov} on
the predictable-branch \texttt{log2\_fold} workload degrades performance to
0.655$\times$ of stock ($p{=}0.002$); the same optimization yields
5.5$\times$ speedup on unpredictable \texttt{cmov\_select}. A fixed kernel
heuristic cannot accommodate both.

\textbf{Key insight.} JIT backend optimization should separate
\emph{mechanism} (the kernel guarantees that a requested native instruction
variant is safe and correct for the current CPU and BPF pattern) from
\emph{policy} (userspace decides which safe variant to use given workload
context, CPU microarchitecture, and deployment constraints). This is
analogous to Halide's algorithm/schedule separation~\cite{halide}, applied
to kernel JIT compilation.

\textbf{Our system.} We design and implement \textsc{BpfReJIT}, a framework
that realizes this separation. The key interface is
\texttt{BPF\_PROG\_JIT\_RECOMPILE}: a new BPF syscall subcommand through
which privileged userspace submits a cryptographically bound policy blob
specifying, per optimization site, which native instruction variant to emit.
The kernel validates each rule against structural constraints and CPU
features, then re-JITs the program without re-verification---the verified
BPF bytecode is never modified.

Five properties distinguish userspace policy from kernel heuristics:
(1)~\emph{updateability} without kernel upgrades;
(2)~\emph{program-level composition} across optimization sites;
(3)~\emph{workload adaptation} based on runtime data distributions;
(4)~\emph{fleet A/B testing} of new policies;
(5)~\emph{ownership} by service teams rather than kernel maintainers.

\textbf{Contributions.}
\begin{itemize}
  \item A characterization of the kernel JIT code quality gap: 0.849$\times$
    exec-time, 0.496$\times$ code-size, with three root causes covering
    89\% of the instruction surplus (\S\ref{sec:motivation}).
  \item Evidence that optimal lowering is policy-sensitive: the same legal
    optimization degrades predictable-branch workloads by 34.5\% while
    improving unpredictable-branch workloads by 5.5$\times$
    (\S\ref{sec:motivation}).
  \item The design and implementation of \textsc{BpfReJIT}: mechanism/policy
    separation for the kernel BPF JIT, with four rule families and a
    declarative extension mechanism (\S\ref{sec:design},
    \S\ref{sec:implementation}).
  \item Evaluation showing ROTATE achieves 1.24$\times$ speedup ($p{=}0.002$),
    60/60 production recompilations succeed, and declarative patterns reduce
    extension cost by 31.8$\times$ (\S\ref{sec:evaluation}).
\end{itemize}

%% ============================================================
\section{Background and Motivation}
\label{sec:motivation}
%% ============================================================

\subsection{The Kernel eBPF JIT}

eBPF programs are compiled from C to BPF bytecode by Clang, then verified by
the kernel's static analysis engine for memory safety and termination. The
verifier produces \emph{extended} (``xlated'') bytecode with resolved helper
calls and rewritten context accesses. The JIT compiler then translates xlated
bytecode to native code via a large \texttt{switch} statement: each BPF
opcode maps to a fixed native instruction sequence. There is no intermediate
representation, no instruction selection pass, and no peephole
optimizer~\cite{linux_bpf_jit}.

Three consequences follow from this design:
(1) multi-byte values that the verifier decomposes into byte loads are
reconstructed with shift-or chains (\texttt{movzbq}+\texttt{shl}+\texttt{or})
instead of single wide loads;
(2) every BPF conditional jump becomes an x86 \texttt{jcc}---there is no
analysis to identify select patterns amenable to \texttt{cmov};
(3) callee-saved registers are unconditionally saved regardless of actual usage
(partially addressed in kernel 7.0-rc2 via \texttt{detect\_reg\_usage()}).

\subsection{Existing BPF Optimizers Cannot Close This Gap}

K2~\cite{k2} (SIGCOMM~'21), Merlin~\cite{merlin} (ASPLOS~'24), and
EPSO~\cite{epso} (ASE~'25) optimize BPF bytecode---the input to the JIT,
not its output. No BPF bytecode rewrite can introduce an x86
\texttt{cmovcc}, fuse shifts into an \texttt{lea}, or replace byte-load
chains with a wide \texttt{mov}: these are native instruction selection
decisions absent from the BPF ISA. The three systems are complementary to
our work, not competing.

\begin{table}[t]
  \centering
  \caption{Comparison with BPF optimization systems.}
  \label{tab:existing_work}
  \small
  \begin{tabular}{lcccc}
    \toprule
    System & Layer & Backend? & Post-load? & Policy ctrl? \\
    \midrule
    K2~\cite{k2}          & BPF bytecode & No  & No  & No  \\
    Merlin~\cite{merlin}  & LLVM IR+BPF  & No  & No  & No  \\
    EPSO~\cite{epso}      & BPF bytecode & No  & No  & No  \\
    \textbf{BpfReJIT}     & \textbf{JIT backend} & \textbf{Yes} & \textbf{Yes} & \textbf{Yes} \\
    \bottomrule
  \end{tabular}
\end{table}

\subsection{Quantifying the Gap}
\label{sec:characterization}

\textbf{Methodology.} We compare the kernel JIT (Linux 6.15.11) against
llvmbpf on 31 pure-JIT microbenchmarks that isolate code generation quality
(no map lookups on the hot path). Each benchmark runs 30 iterations of 1{,}000
repeats; we report the per-benchmark median of per-iteration medians. We use
paired Wilcoxon signed-rank tests with Benjamini-Hochberg correction, and
bootstrap 95\% CIs for suite-level geomeans. The platform is an Intel Core
Ultra 9 285K (Arrow Lake-S), performance governor, turbo disabled.

\begin{table}[t]
  \centering
  \caption{Kernel JIT code quality gap: suite-level summary.}
  \label{tab:exec_summary}
  \begin{tabular}{lrrr}
    \toprule
    Suite & Geomean (L/K) & 95\% CI & Sig.\ / Total \\
    \midrule
    Pure-JIT (31 auth.) & 0.849$\times$ & [0.834, 0.865] & 25/31 \\
    Runtime (9)         & 0.829$\times$ & [0.801, 0.859] & 7/9   \\
    \bottomrule
  \end{tabular}
\end{table}

\begin{table}[t]
  \centering
  \caption{Selected per-benchmark results (31 authoritative).}
  \label{tab:per_bench}
  \small
  \begin{tabular}{lrrrr}
    \toprule
    Benchmark & llvmbpf & Kernel & Exec L/K & Code L/K \\
    \midrule
    \texttt{branch\_layout}    & 162\,ns  & 561\,ns  & \textbf{0.289$\times$} & 0.354$\times$ \\
    \texttt{stride\_load\_16}  & 104\,ns  & 246\,ns  & 0.423$\times$          & 0.332$\times$ \\
    \texttt{binary\_search}    & 213\,ns  & 442\,ns  & 0.481$\times$          & 0.449$\times$ \\
    \texttt{switch\_dispatch}  & 212\,ns  & 268\,ns  & 0.789$\times$          & 0.800$\times$ \\
    \texttt{log2\_fold}        & 312\,ns  & 311\,ns  & 1.006$\times$          & 0.659$\times$ \\
    \texttt{bitcount}          & 4647\,ns & 2985\,ns & 1.556$\times$          & 0.446$\times$ \\
    \texttt{code\_clone\_8}    & 2267\,ns & 1207\,ns & \textbf{1.878$\times$} & 0.542$\times$ \\
    \midrule
    \multicolumn{3}{l}{\textit{Geomean (31 benchmarks)}} & 0.849$\times$ & 0.496$\times$ \\
    \bottomrule
  \end{tabular}
\end{table}

Table~\ref{tab:exec_summary} summarizes the gap. llvmbpf is faster on
21/31 benchmarks; the kernel wins on 10/31, primarily loop-dominated
workloads where the JIT's simpler code has lower I-cache pressure.
Table~\ref{tab:per_bench} shows selected results spanning the 6.5$\times$
range (0.289$\times$ to 1.878$\times$). Code size tells a cleaner story:
llvmbpf is smaller on \emph{all} 31 benchmarks (geomean 0.496$\times$),
confirming the gap is systematic.

\textbf{Root causes.} JIT-dump analysis of 22 benchmarks decomposes the
kernel's instruction surplus:

\begin{table}[t]
  \centering
  \caption{Sources of kernel JIT instruction surplus.}
  \label{tab:surplus}
  \begin{tabular}{lrl}
    \toprule
    Source & Share & Mechanism \\
    \midrule
    Byte-load recomposition & 50.7\% & \texttt{movzbq+shl+or} chains \\
    Absent conditional moves& 19.9\% & No cmov; 31 vs.\ 0 \\
    Fixed callee-saved regs & 18.5\% & Unconditional \texttt{push/pop} \\
    Other                   & 10.9\% & NOPs, zero-ext, misc \\
    \bottomrule
  \end{tabular}
\end{table}

These three mechanisms account for 89\% of the surplus. The isolated
time penalty for byte-load recomposition is 2.24$\times$.

\subsection{The Policy-Sensitivity Problem}
\label{sec:policy_sensitivity}

A fixed kernel heuristic would address the gap only if every optimization
were unconditionally profitable. Our ablation study shows otherwise:

\begin{table}[t]
  \centering
  \caption{Cmov ablation in llvmbpf: profitability depends on workload.}
  \label{tab:cmov_ablation}
  \begin{tabular}{lrrr}
    \toprule
    Benchmark & Normal & No-cmov & $\Delta$ \\
    \midrule
    \texttt{switch\_dispatch}  & 213\,ns & 269\,ns & \textbf{+26.3\%} \\
    \texttt{binary\_search}    & 205\,ns & 229\,ns & \textbf{+11.7\%} \\
    \texttt{bounds\_ladder}    & 83\,ns  & 68\,ns  & $-18.1\%$ \\
    \texttt{large\_mixed\_500} & 415\,ns & 316\,ns & $-23.9\%$ \\
    \bottomrule
  \end{tabular}
\end{table}

Disabling \texttt{cmov} \emph{hurts} programs with unpredictable branches
(\texttt{switch\_dispatch}: data-dependent dispatch) but \emph{helps}
programs with predictable branches (\texttt{bounds\_ladder}: monotone
thresholds). The same legal lowering helps one workload class and hurts
another by comparable magnitudes. A fixed ``always-cmov'' kernel heuristic
would degrade the second class; ``never-cmov'' would degrade the first.

The mechanism we need is not a better heuristic but a better
\emph{interface}: the kernel should enforce which lowerings are \emph{safe},
while the operator specifies which safe lowerings are \emph{profitable}.

%% ============================================================
\section{Design}
\label{sec:design}
%% ============================================================

\textsc{BpfReJIT} is organized around a strict \emph{mechanism/policy
separation}. The kernel \emph{legality plane} enforces which native
instruction choices are safe for a given BPF pattern on the current CPU.
The userspace \emph{policy plane} decides which safe choices to request
based on workload profile, CPU features, and deployment constraints.

\subsection{Architecture Overview}

Figure~\ref{fig:architecture} shows the two-plane architecture.
The interface between planes is \texttt{BPF\_PROG\_JIT\_RECOMPILE}, a new
BPF syscall subcommand. Userspace submits a \emph{policy blob}---a compact
binary specification of per-site native instruction choices---attached to a
sealed memfd. The kernel validates and applies the blob, re-JITting the
program without re-verification.

\begin{figure}[t]
  \centering
  \small
  \begin{tabular}{|p{0.42\linewidth}|}
    \hline
    \textbf{Userspace Policy Plane} \\
    \hline
    Pattern scanning on xlated bytecode \\
    CPU + workload policy selection \\
    Fleet-level A/B rollout \\
    Code-size budget management \\
    Policy blob generation \\
    \hline
    \multicolumn{1}{c}{$\downarrow$ \texttt{BPF\_PROG\_JIT\_RECOMPILE}} \\
    \hline
    \textbf{Kernel Legality Plane} \\
    \hline
    Digest binding (prog\_tag) \\
    Pattern matching + variable binding \\
    CPU feature gating \\
    Structural constraint validation \\
    Canonical-parameter extraction \\
    Parameterized arch-specific emitters \\
    Fail-closed fallback per rule \\
    \hline
  \end{tabular}
  \caption{Two-plane architecture of \textsc{BpfReJIT}. Userspace controls
    \emph{what} to optimize; the kernel controls \emph{how} safely.}
  \label{fig:architecture}
\end{figure}

\textbf{Core invariant.} The verified BPF bytecode (\texttt{xlated\_prog\_len})
is never modified. Policy affects only native code emission. The verifier's
safety guarantees remain intact.

\textbf{Fail-closed design.} If the policy blob's digest does not match the
loaded program, the entire blob is rejected. If an individual rule's pattern
match or constraint check fails, that site falls back silently to stock JIT
emission. A malformed policy cannot crash or incorrectly optimize a program.

\subsection{The BPF\_PROG\_JIT\_RECOMPILE Interface}

\begin{lstlisting}[language=C,caption=Core syscall interface.]
union bpf_attr {
    struct { /* BPF_PROG_JIT_RECOMPILE */
        __u32 prog_fd;    /* loaded program */
        __u32 policy_fd;  /* sealed memfd; 0=revert */
        __u32 flags;
        __u32 log_level;
        __aligned_u64 log_buf;
        __u32 log_size;
    } jit_recompile;
};
\end{lstlisting}

When \texttt{policy\_fd == 0}, the kernel reverts to stock JIT, giving
operators a simple undo. The interface enables an iterative optimization
loop: load $\to$ run $\to$ profile $\to$ generate policy $\to$ re-JIT $\to$
measure $\to$ refine---all without program restart or re-verification.

\subsection{Rule Families}

We define four rule families covering the root causes from
\S\ref{sec:characterization}:

\begin{table}[t]
  \centering
  \caption{Rule families in \textsc{BpfReJIT}.}
  \label{tab:rule_families}
  \small
  \begin{tabular}{llp{3.8cm}}
    \toprule
    Family & Choices & Addresses \\
    \midrule
    \textsc{Cond\_Select} & \texttt{cmovcc} vs.\ branch & Absent cmov (policy-sensitive) \\
    \textsc{Wide\_Mem}    & Wide load vs.\ byte ladder   & Byte-load recomposition \\
    \textsc{Rotate}       & \texttt{rorx}/\texttt{ror} vs.\ shifts & Hash/checksum rotate idioms \\
    \textsc{Addr\_Calc}   & \texttt{lea} vs.\ shift+add  & Address computation fusion \\
    \bottomrule
  \end{tabular}
\end{table}

Each family defines: (1) the BPF instruction pattern it matches,
(2) the set of legal native instruction variants, (3) structural
constraints (straight-line, no side effects, no interior branch edges),
and (4) CPU feature requirements (e.g., BMI2 for \texttt{rorx}).

\subsection{Security Model}

\textbf{No native code upload.} Userspace specifies \emph{which named
variant} to emit (e.g., \texttt{BPF\_JIT\_SEL\_CMOVCC}), not machine bytes.
Emission is performed by kernel emitter functions.

\textbf{Digest binding.} The policy blob carries \texttt{prog\_tag}
(SHA-256 of xlated bytecode). Tag mismatch rejects the entire blob.
The blob is passed via a sealed memfd (\texttt{F\_SEAL\_WRITE} +
\texttt{F\_SEAL\_SHRINK}), preventing TOCTOU attacks.

\textbf{CPU feature gating.} Each rule declares required CPU features;
the kernel checks CPUID before applying. A \texttt{rorx} request on a
CPU without BMI2 is silently rejected.

\textbf{Privilege.} \texttt{BPF\_PROG\_JIT\_RECOMPILE} requires
\texttt{CAP\_BPF}, the same privilege as \texttt{BPF\_PROG\_LOAD}.

\subsection{Declarative Pattern Extension (v5)}
\label{sec:v5}

The initial v4 design uses a closed \texttt{rule\_kind} enum: each
optimization has a hardcoded kernel validator. Adding the final ROTATE
support required 1{,}018 LOC (670 kernel + 348 userspace).

v5 moves pattern descriptions into the policy blob. The kernel retains a
small fixed set of \emph{canonical forms} (e.g., \textsc{CF\_Rotate},
\textsc{CF\_Select}), each with a parameterized emitter. A new pattern
variant that maps to an existing canonical form requires only a userspace
descriptor edit---no kernel changes. The pipeline is:

\begin{center}
\small
pattern descriptor $\to$ variable bindings $\to$ canonical parameters $\to$ parameterized emitter
\end{center}

This achieves a 31.8$\times$ reduction in extension cost for same-form
variants (32 LOC vs.\ 1{,}018 LOC; \S\ref{sec:eval_extensibility}).

%% ============================================================
\section{Implementation}
\label{sec:implementation}
%% ============================================================

\textsc{BpfReJIT} modifies the Linux kernel (based on v7.0-rc2) in
approximately 1{,}800 LOC across 7 files, plus a shared userspace
scanner/blob-builder.

\textbf{Kernel components.}
The implementation adds: (1)~\texttt{BPF\_PROG\_JIT\_RECOMPILE} plumbing
and policy blob parsing (${\sim}$350 LOC); (2)~a generic per-rule validator
checking structural constraints, pattern matching, and CPU features
(${\sim}$150 LOC); (3)~four canonical-form validators and x86 emitters
(${\sim}$640 LOC total); (4)~rule dispatch in the JIT main loop
(${\sim}$120 LOC); (5)~re-JIT image installation (${\sim}$80 LOC; the
current prototype stops short of fully RCU-protected hot-swap).

\textbf{Rule dispatch.} We instrument \texttt{do\_jit()} to check for
applicable rules before each BPF instruction:

\begin{lstlisting}[language=C,caption=Rule dispatch in do\_jit().]
for (i = 1; i <= insn_cnt; i++, insn++) {
    rule = bpf_jit_lookup_rule(ruleset, i-1);
    if (rule) {
        int len = bpf_jit_emit_rule_x86(
            prog, &ctx, rule, addrs, i-1);
        if (len >= 0) {
            i    += rule->site_len - 1;
            insn += rule->site_len - 1;
            continue;
        }
        /* fallback to stock emission */
    }
    /* stock switch-case ... */
}
\end{lstlisting}

Rule lookup is O(1) via sorted array indexing. Failed emission falls
through to stock processing without side effects.

\textbf{Validation layers.} Each rule passes through four validation layers:
(1)~structural (bounds, subprogram containment, known \texttt{rule\_kind},
CPU features); (2)~pattern (straight-line, no helpers/stores/atomics, opcode
matching, variable binding); (3)~canonical-parameter extraction (binding
table maps matched variables to emitter parameters); (4)~form-specific
checks (${\sim}$50--100 LOC per form; e.g., ROTATE verifies shift amounts
sum to operand width).

\textbf{Userspace scanner.} The scanner retrieves xlated bytecode via
\texttt{bpf\_prog\_get\_info\_by\_fd()}, applies declarative pattern
descriptors for each family, extracts site features, applies CPU/workload
policy, serializes v5 rules into a policy blob, and invokes
\texttt{BPF\_PROG\_JIT\_RECOMPILE} via a sealed memfd.

%% ============================================================
\section{Evaluation}
\label{sec:evaluation}
%% ============================================================

We organize evaluation around five research questions.

\textbf{Setup.} Microbenchmark re-JIT experiments use a rigorous
configuration: 5 warmups, 10 measured iterations, 1{,}000 repeats per
iteration, single guest vCPU pinned to host CPU~23, host governor
\texttt{performance}, turbo disabled. The guest runs our modified
kernel~7.0-rc2 inside QEMU/KVM. We report median \texttt{exec\_ns} and
paired Wilcoxon signed-rank $p$-values with BH correction.

\subsection{RQ1: How Large Is the Code Quality Gap?}

The characterization in \S\ref{sec:characterization} establishes the gap.
The kernel JIT produces code that is 15\% slower (geomean 0.849$\times$)
and 2$\times$ larger (geomean 0.496$\times$) than LLVM~-O3 on identical
BPF programs. Three backend deficiencies account for 89\% of the
instruction surplus (Table~\ref{tab:surplus}). The gap is confirmed on
36 unique real-world programs from Cilium and libbpf-bootstrap: code-size
geomean 0.618$\times$, exec-time geomean 0.514$\times$ (14 unique runnable).

\textbf{Takeaway.} The kernel JIT leaves substantial performance on the
table, and the root causes are in the JIT backend, not the BPF bytecode.

\subsection{RQ2: Can a Fixed Kernel Policy Close the Gap?}

No. Table~\ref{tab:cmov_ablation} shows that \texttt{cmov} profitability
varies by ${\pm}$26\% across workloads. Our rigorous re-JIT experiment
makes this concrete: forcing CMOV on the predictable-branch
\texttt{log2\_fold} degrades performance from 383.5\,ns to 590.0\,ns
(0.655$\times$, $p{=}0.002$, stock wins all 10 iterations), while the
characterization shows that programs \emph{without} \texttt{cmov}
(\texttt{switch\_dispatch}) suffer 26\% overhead from the same absence.

\begin{table*}[t]
  \centering
  \caption{Per-family re-JIT results (rigorous pinned 10$\times$1000 configuration).
    Speedup = stock median / re-JIT median; values $<$1 indicate policy hurts.}
  \label{tab:per_family}
  \small
  \begin{tabular}{llrrrrrl}
    \toprule
    Policy & Benchmark & Stock (ns) & Re-JIT (ns) & Speedup & Code $\Delta$ & $p$ & Interpretation \\
    \midrule
    \textsc{Cond\_Select}  & \texttt{log2\_fold}       & 383.5 & 590.0 & 0.655$\times$ & +5.3\% & 0.002 & Policy-sensitivity evidence \\
    \textsc{Rotate}+\textsc{Wide} & \texttt{rotate64\_hash} & 68.0  & 57.0  & 1.193$\times$ & $-$35.0\% & 0.002 & Combined optimization \\
    \textsc{Rotate}        & \texttt{pkt\_rss\_hash}   & 21.0  & 17.0  & 1.235$\times$ & $-$15.3\% & 0.002 & Pure rotate gain \\
    \textsc{Addr\_Calc}    & \texttt{stride\_load\_16} & 293.0 & 281.0 & 1.052$\times$ & $-$1.2\%  & 0.012 & Modest but significant \\
    \textsc{Wide\_Mem}     & \texttt{load\_byte\_recomp} & 223.0 & 226.0 & 0.978$\times$ & $-$2.8\% & 0.164 & Code-size win only \\
    \bottomrule
  \end{tabular}
\end{table*}

A single fixed policy cannot realize the ROTATE gains
(\texttt{packet\_rss\_hash}: 1.235$\times$) without also inflicting the
CMOV regression (\texttt{log2\_fold}: 0.655$\times$). The correct
interface is a policy plane controlled by the operator.

\textbf{Takeaway.} Fixed kernel heuristics are falsified: the same legal
optimization produces a 34.5\% degradation or a 5.5$\times$ improvement
depending on branch predictability.

\subsection{RQ3: Can BpfReJIT Close the Gap Safely?}

We validate safety through three properties.

\textbf{Bytecode invariant.} Re-JIT preserves the verified bytecode exactly:
\texttt{log2\_fold} xlated bytes remain 1{,}112 before and after CMOV
recompile; \texttt{rotate64\_hash} remains 7{,}984 while JIT bytes shrink
from 3{,}559 to 2{,}313.

\textbf{Production recompilation.} We exercise the v5 recompile path on
two production-oriented corpora:
\begin{itemize}
  \item \textbf{Expanded corpus} (Calico, Suricata): 13 objects, 61
    perf-capable programs, 60 baseline compiles, \textbf{60/60 v5 recompiles
    succeed}, 45 with optimization sites applied, zero syscall failures.
  \item \textbf{Original corpus} (Katran, xdp-tools, xdp-tutorial): 38
    programs compile, 15 with sites, \textbf{15/15 recompiles succeed}.
\end{itemize}

The framework detects 197 CMOV + 108 WIDE + 10 ROTATE sites across the
expanded corpus programs. Code-size changes are small (geomean
1.001$\times$), with the largest reduction being \texttt{calico\_tc\_maglev}
at $-$4.4\% (10 ROTATE sites).

\textbf{Zero overhead.} Programs with no matching sites incur no runtime
or code-size overhead. Rule lookup in \texttt{do\_jit()} returns null
immediately; programs that never receive a recompile call follow the
identical code path as stock JIT.

\textbf{Takeaway.} \textsc{BpfReJIT} preserves the bytecode invariant,
achieves 100\% recompile success on production programs, and imposes
zero overhead on non-targeted programs.

\subsection{RQ4: How Effective Is Each Rule Family?}

Table~\ref{tab:per_family} presents per-family results.

\textbf{ROTATE} delivers the strongest gains: \texttt{packet\_rss\_hash}
improves from 21.0\,ns to 17.0\,ns (1.235$\times$, $p{=}0.002$) with
15.3\% code reduction. \texttt{rotate64\_hash} under the combined
ROTATE+WIDE policy improves from 68.0\,ns to 57.0\,ns (1.193$\times$,
$p{=}0.002$) with 35.0\% code reduction. These gains are robust: re-JIT
wins all 10 measured iterations in both cases.

\textbf{ADDR\_CALC} gives a modest but significant 5.2\% speedup on
\texttt{stride\_load\_16} ($p{=}0.012$), replacing multi-instruction
shift+add sequences with single \texttt{lea} instructions.

\textbf{WIDE\_MEM} shows code-size reduction ($-$2.8\%) but the runtime
difference is not statistically significant ($p{=}0.164$). We report this
as a limitation: the current evidence for WIDE\_MEM is stronger for code
size than execution time.

\textbf{COND\_SELECT} demonstrates the policy-sensitivity thesis rather than
a positive gain: forcing CMOV on \texttt{log2\_fold} produces a 34.5\%
\emph{regression}, validating that this family must be policy-controlled.
The characterization ablation (Table~\ref{tab:cmov_ablation}) confirms the
opposite direction: removing CMOV hurts \texttt{switch\_dispatch} by 26.3\%.

\textbf{Takeaway.} Two families deliver statistically significant speedups;
one delivers code-size wins; one demonstrates that unconditional application
hurts---exactly the distribution that motivates mechanism/policy separation.

\subsection{RQ5: How Extensible Is the Framework?}
\label{sec:eval_extensibility}

Table~\ref{tab:extensibility} quantifies extension cost.

\begin{table}[t]
  \centering
  \caption{Extension cost: v4 closed-enum vs.\ v5 declarative patterns.}
  \label{tab:extensibility}
  \small
  \begin{tabular}{lrrr}
    \toprule
    Scenario & Kernel & Userspace & Total \\
    \midrule
    v4: add ROTATE family          & 670 & 348 & 1{,}018 \\
    v5: new variant, same form     & 0   & 32  & 32      \\
    v5: new canonical form         & 111 & 55  & 166     \\
    \bottomrule
  \end{tabular}
\end{table}

Under v4, adding ROTATE required 670 kernel LOC (UAPI definitions,
validators, emitters) plus 348 userspace LOC (scanners, CLI plumbing)
across 3 kernel files and multiple userspace files.

Under v5, adding a new ROTATE \emph{variant} that maps to the existing
\textsc{CF\_Rotate} canonical form requires \textbf{0 kernel LOC} and
\textbf{32 userspace LOC}---a single pattern descriptor with constraint
list and canonical bindings. Even a brand-new canonical form
(e.g., \textsc{CF\_Addr\_Calc}) costs only 166 LOC because the pattern
matching and binding machinery are shared.

\textbf{Takeaway.} Declarative patterns reduce same-form extension cost by
31.8$\times$ (1{,}018 $\to$ 32 LOC), with zero kernel files touched.

%% ============================================================
\section{Discussion}
\label{sec:discussion}
%% ============================================================

\textbf{Falsification condition.} If fixed kernel heuristics recovered the
same gains on all tested workloads, the correct conclusion would be to add
kernel peepholes. Our evidence rejects this: always-CMOV degrades
\texttt{log2\_fold} by 34.5\% ($p{=}0.002$). The remaining gap to a fully
rigorous falsification is cross-microarchitecture validation (ARM64, older
Intel). We conjecture that policy sensitivity increases on architectures
with different branch predictor designs.

\textbf{LLVM as upper bound.} llvmbpf serves as an upper bound reference,
not a direct competitor. Running LLVM -O3 in the kernel is impractical;
\textsc{BpfReJIT} captures a portion of the available headroom in a
kernel-compatible way.

\textbf{Complementarity.} K2, Merlin, and EPSO optimize BPF bytecode
\emph{before} JIT; \textsc{BpfReJIT} optimizes native code \emph{after}.
A deployment could use both layers without interference.

\textbf{Limitations.} Four limitations bound our current results.
(1)~\textsc{Wide\_Mem} does not yet show a statistically significant latency
win ($p{=}0.164$). (2)~Cilium programs fail to load in our test environment
due to BTF/section compatibility issues. (3)~The re-JIT installation path
is not yet fully RCU-protected for concurrent execution. (4)~The v5 pattern
matcher uses exact opcode comparison rather than masked matching, causing
descriptor explosion for families with many opcode variants.

%% ============================================================
\section{Related Work}
\label{sec:related}
%% ============================================================

\textbf{BPF-level optimization.}
K2~\cite{k2} applies stochastic superoptimization to BPF bytecode.
Merlin~\cite{merlin} lifts BPF to LLVM IR for optimization.
EPSO~\cite{epso} uses evolutionary synthesis at the bytecode level.
All three operate above the JIT and cannot influence native instruction
selection. \textsc{BpfReJIT} differs on three axes: it operates at the JIT
backend layer, supports post-load re-optimization without re-verification,
and enables per-site selective application to handle policy sensitivity.
The systems are complementary: bytecode optimizers improve the JIT's input;
\textsc{BpfReJIT} improves its output.

\textbf{Superoptimization and peephole optimization.}
Massalin's superoptimizer~\cite{superoptimizer} and
STOKE~\cite{stoke} find optimal native sequences.
Bansal and Aiken~\cite{credal} automate peephole rule generation.
Our problem is different: the kernel already compiles correctly; the
question is which \emph{legal native variant} to emit, not whether to
replace a sequence.

\textbf{Feedback-directed and JIT compilation.}
Profile-guided optimization~\cite{feedback_directed} and adaptive JIT
techniques~\cite{jit_survey} are well-established. Our contribution
applies this principle to the kernel BPF JIT specifically: post-load
re-JIT enables production programs to be re-optimized based on actual
workload behavior without restart.

\textbf{Compiler backend design.}
LLVM's TableGen-based instruction selection~\cite{llvm} describes
target-specific lowering patterns. Our canonical forms can be seen as a
safety-restricted subset of instruction selection, constrained to
pre-verified equivalence-preserving lowerings.

\textbf{Algorithm/schedule separation.}
Halide~\cite{halide} separates computation from scheduling, enabling
independent optimization of each. \textsc{BpfReJIT} applies the same
principle: BPF bytecode specifies the algorithm; the policy blob
specifies the native implementation strategy.

%% ============================================================
\section{Conclusion}
\label{sec:conclusion}
%% ============================================================

We characterized the performance gap between the Linux kernel eBPF JIT and
LLVM~-O3: a 15\% execution-time gap and 2$\times$ code-size gap rooted in
three backend lowering deficiencies. We showed that these deficiencies
cannot be fixed by kernel heuristics alone because optimal lowering is
policy-sensitive---the same optimization degrades predictable-branch
workloads by 34.5\% while improving unpredictable-branch workloads.

\textsc{BpfReJIT} separates mechanism (kernel safety) from policy (userspace
profitability) via \texttt{BPF\_PROG\_JIT\_RECOMPILE}. Four rule families
cover the identified gaps. Evaluation shows statistically significant speedups
for ROTATE (1.24$\times$, $p{=}0.002$), 100\% recompile success on 60
production programs, and a 31.8$\times$ reduction in extension cost through
declarative patterns. The framework makes explicit a design choice that
has been implicit in the kernel JIT: backend optimization decisions are
policy decisions, and policy should be externalized.

\begin{acks}
Anonymized for submission.
\end{acks}

\bibliographystyle{ACM-Reference-Format}
\bibliography{references}

\appendix

\section{Benchmark Descriptions}
\label{appendix:benchmarks}

\begin{table}[h]
  \centering
  \caption{Selected benchmark descriptions (31 authoritative suite).}
  \label{tab:bench_desc}
  \small
  \begin{tabular}{llp{4.5cm}}
    \toprule
    Benchmark & Category & Description \\
    \midrule
    \texttt{branch\_layout}        & ctrl-flow & Conditional branches; tests layout opt \\
    \texttt{switch\_dispatch}      & ctrl-flow & Multi-way dispatch; high cmov density in LLVM \\
    \texttt{binary\_search}        & ctrl-flow & Binary search; variable branch predictability \\
    \texttt{load\_byte\_recompose} & memory    & 8-byte recompose from byte loads \\
    \texttt{stride\_load\_16}      & memory    & 16-byte strided loads; LEA candidates \\
    \texttt{log2\_fold}            & ALU       & Integer log2; predictable branches \\
    \texttt{cmov\_select}          & ctrl-flow & Narrow select diamond; cmov isolation \\
    \texttt{rotate64\_hash}        & ALU       & 64-bit rotate hash; rotate fusion \\
    \texttt{packet\_rss\_hash}     & hash      & RSS hash computation; 11 rotate sites \\
    \bottomrule
  \end{tabular}
\end{table}

\section{Policy Blob Binary Format}
\label{appendix:blob}

The policy blob is a flat binary structure:
(1)~\texttt{bpf\_jit\_policy\_hdr} (32 bytes) containing magic, version,
rule count, program instruction count, \texttt{prog\_tag} digest, and
architecture ID;
(2)~array of \texttt{bpf\_jit\_rewrite\_rule} (variable-length in v5) each
specifying site offset, site length, rule kind, native choice, CPU feature
requirements, and (in v5) inline pattern descriptors with constraint arrays
and canonical binding tables;
(3)~rule-specific payload data referenced by offset/length fields.

The blob is sealed via \texttt{F\_SEAL\_WRITE} and \texttt{F\_SEAL\_SHRINK}
before submission, preventing TOCTOU modification between digest validation
and rule application.

\section{Production Corpus Details}
\label{appendix:corpus}

\begin{table}[h]
  \centering
  \caption{Directive site census across production corpus.}
  \label{tab:census}
  \small
  \begin{tabular}{lrrrrr}
    \toprule
    Project & Objects & CMOV & WIDE & ROTATE & Total \\
    \midrule
    Calico   & 8 & 196 & 108   & 10 & 314  \\
    loxilb   & 3 & 98  & 1{,}729 & 0  & 1{,}827\\
    Suricata & 2 & 2   & 0     & 0  & 2    \\
    Katran   & 5 & 6   & 4     & 0  & 10   \\
    xdp-tools& 4 & 4   & 2     & 0  & 6    \\
    \midrule
    \textbf{Total} & \textbf{22} & \textbf{306} & \textbf{1{,}843} & \textbf{10} & \textbf{2{,}159} \\
    \bottomrule
  \end{tabular}
\end{table}

Note: loxilb objects (1{,}827 sites) could not be loaded through libbpf
due to non-standard ELF section names; these counts are inventory-only,
not validated via recompile. Cilium objects (3 objects, 658 CMOV sites
from scanner fallback) failed at \texttt{bpf\_object\_\_open\_file()}.

\end{document}
